% !TeX program = xelatex
% UTF-8. Overleaf: choose XeLaTeX; this file is self-contained.
% Run twice for the table of contents.
% Options for packages loaded elsewhere
\PassOptionsToPackage{unicode}{hyperref}
\PassOptionsToPackage{hyphens}{url}
\PassOptionsToPackage{dvipsnames,svgnames,x11names}{xcolor}
\documentclass[
  11pt,
  a4paper,
]{article}
\usepackage{xcolor}
\usepackage[margin=20mm]{geometry}
\usepackage{amsmath,amssymb}
\setcounter{secnumdepth}{-\maxdimen} % remove section numbering
\usepackage{iftex}
\ifPDFTeX
  \usepackage[T1]{fontenc}
  \usepackage[utf8]{inputenc}
  \usepackage{textcomp} % provide euro and other symbols
\else % if luatex or xetex
  \usepackage{unicode-math} % this also loads fontspec
  \defaultfontfeatures{Scale=MatchLowercase}
  \defaultfontfeatures[\rmfamily]{Ligatures=TeX,Scale=1}
\fi
\usepackage{lmodern}
\ifPDFTeX\else
  % xetex/luatex font selection
  \setmainfont[BoldFont=texgyretermes-bold.otf,ItalicFont=texgyretermes-italic.otf,BoldItalicFont=texgyretermes-bolditalic.otf]{texgyretermes-regular.otf}
  \setsansfont[BoldFont=texgyreheros-bold.otf,ItalicFont=texgyreheros-italic.otf,BoldItalicFont=texgyreheros-bolditalic.otf]{texgyreheros-regular.otf}
  \setmonofont[BoldFont=texgyrecursor-bold.otf,ItalicFont=texgyrecursor-italic.otf,BoldItalicFont=texgyrecursor-bolditalic.otf]{texgyrecursor-regular.otf}
\fi
% Use upquote if available, for straight quotes in verbatim environments
\IfFileExists{upquote.sty}{\usepackage{upquote}}{}
\IfFileExists{microtype.sty}{% use microtype if available
  \usepackage[]{microtype}
  \UseMicrotypeSet[protrusion]{basicmath} % disable protrusion for tt fonts
}{}
\makeatletter
\@ifundefined{KOMAClassName}{% if non-KOMA class
  \IfFileExists{parskip.sty}{%
    \usepackage{parskip}
  }{% else
    \setlength{\parindent}{0pt}
    \setlength{\parskip}{6pt plus 2pt minus 1pt}}
}{% if KOMA class
  \KOMAoptions{parskip=half}}
\makeatother
\usepackage{color}
\usepackage{fancyvrb}
\newcommand{\VerbBar}{|}
\newcommand{\VERB}{\Verb[commandchars=\\\{\}]}
\DefineVerbatimEnvironment{Highlighting}{Verbatim}{commandchars=\\\{\}}
% Add ',fontsize=\small' for more characters per line
\newenvironment{Shaded}{}{}
\newcommand{\AlertTok}[1]{\textcolor[rgb]{1.00,0.00,0.00}{\textbf{#1}}}
\newcommand{\AnnotationTok}[1]{\textcolor[rgb]{0.38,0.63,0.69}{\textbf{\textit{#1}}}}
\newcommand{\AttributeTok}[1]{\textcolor[rgb]{0.49,0.56,0.16}{#1}}
\newcommand{\BaseNTok}[1]{\textcolor[rgb]{0.25,0.63,0.44}{#1}}
\newcommand{\BuiltInTok}[1]{\textcolor[rgb]{0.00,0.50,0.00}{#1}}
\newcommand{\CharTok}[1]{\textcolor[rgb]{0.25,0.44,0.63}{#1}}
\newcommand{\CommentTok}[1]{\textcolor[rgb]{0.38,0.63,0.69}{\textit{#1}}}
\newcommand{\CommentVarTok}[1]{\textcolor[rgb]{0.38,0.63,0.69}{\textbf{\textit{#1}}}}
\newcommand{\ConstantTok}[1]{\textcolor[rgb]{0.53,0.00,0.00}{#1}}
\newcommand{\ControlFlowTok}[1]{\textcolor[rgb]{0.00,0.44,0.13}{\textbf{#1}}}
\newcommand{\DataTypeTok}[1]{\textcolor[rgb]{0.56,0.13,0.00}{#1}}
\newcommand{\DecValTok}[1]{\textcolor[rgb]{0.25,0.63,0.44}{#1}}
\newcommand{\DocumentationTok}[1]{\textcolor[rgb]{0.73,0.13,0.13}{\textit{#1}}}
\newcommand{\ErrorTok}[1]{\textcolor[rgb]{1.00,0.00,0.00}{\textbf{#1}}}
\newcommand{\ExtensionTok}[1]{#1}
\newcommand{\FloatTok}[1]{\textcolor[rgb]{0.25,0.63,0.44}{#1}}
\newcommand{\FunctionTok}[1]{\textcolor[rgb]{0.02,0.16,0.49}{#1}}
\newcommand{\ImportTok}[1]{\textcolor[rgb]{0.00,0.50,0.00}{\textbf{#1}}}
\newcommand{\InformationTok}[1]{\textcolor[rgb]{0.38,0.63,0.69}{\textbf{\textit{#1}}}}
\newcommand{\KeywordTok}[1]{\textcolor[rgb]{0.00,0.44,0.13}{\textbf{#1}}}
\newcommand{\NormalTok}[1]{#1}
\newcommand{\OperatorTok}[1]{\textcolor[rgb]{0.40,0.40,0.40}{#1}}
\newcommand{\OtherTok}[1]{\textcolor[rgb]{0.00,0.44,0.13}{#1}}
\newcommand{\PreprocessorTok}[1]{\textcolor[rgb]{0.74,0.48,0.00}{#1}}
\newcommand{\RegionMarkerTok}[1]{#1}
\newcommand{\SpecialCharTok}[1]{\textcolor[rgb]{0.25,0.44,0.63}{#1}}
\newcommand{\SpecialStringTok}[1]{\textcolor[rgb]{0.73,0.40,0.53}{#1}}
\newcommand{\StringTok}[1]{\textcolor[rgb]{0.25,0.44,0.63}{#1}}
\newcommand{\VariableTok}[1]{\textcolor[rgb]{0.10,0.09,0.49}{#1}}
\newcommand{\VerbatimStringTok}[1]{\textcolor[rgb]{0.25,0.44,0.63}{#1}}
\newcommand{\WarningTok}[1]{\textcolor[rgb]{0.38,0.63,0.69}{\textbf{\textit{#1}}}}
\usepackage{longtable,booktabs,array}
\usepackage{caption}
\captionsetup[table]{skip=6pt}
\newcounter{none} % for unnumbered tables
\usepackage{calc} % for calculating minipage widths
% Correct order of tables after \paragraph or \subparagraph
\usepackage{etoolbox}
\makeatletter
\patchcmd\longtable{\par}{\if@noskipsec\mbox{}\fi\par}{}{}
\makeatother
\setlength{\emergencystretch}{3em} % prevent overfull lines
\providecommand{\tightlist}{%
  \setlength{\itemsep}{0pt}\setlength{\parskip}{0pt}}
% Portable Vietnamese Unicode document; compile with XeLaTeX or LuaLaTeX.
\usepackage{fvextra}
\DefineVerbatimEnvironment{verbatim}{Verbatim}{fontsize=\footnotesize,breaklines=true}
\usepackage{etoolbox}
\usepackage{newunicodechar}
\newunicodechar{↔}{\ensuremath{\leftrightarrow}}
\AtBeginEnvironment{longtable}{\footnotesize\raggedright}
\setlength{\tabcolsep}{4pt}
\renewcommand{\arraystretch}{1.22}
\setlength{\emergencystretch}{3em}
\setcounter{tocdepth}{2}
\renewcommand{\contentsname}{Mục lục}
\AtBeginDocument{\hypersetup{pdftitle={Quỷ Mù REV P - Hướng dẫn dựng vách và làm cơ chế cửa}}}
\usepackage{bookmark}
\IfFileExists{xurl.sty}{\usepackage{xurl}}{} % add URL line breaks if available
\urlstyle{same}
% fallback for those not using the hyperref driver hyperxmp:
\makeatletter
\@ifundefined{xmpquote}{\newcommand{\xmpquote}[1]{#1}}{}
\makeatother
\hypersetup{
  pdftitle={Quỷ Mù REV P},
  colorlinks=true,
  linkcolor={teal},
  filecolor={Maroon},
  citecolor={Blue},
  urlcolor={teal},
  pdfcreator={LaTeX via pandoc}}

\title{Quỷ Mù REV P}
\usepackage{etoolbox}
\makeatletter
\providecommand{\subtitle}[1]{% add subtitle to \maketitle
  \apptocmd{\@title}{\par {\large #1 \par}}{}{}
}
\makeatother
\subtitle{Hướng dẫn dựng vách và làm cơ chế cửa}
\author{}
\date{17/10/2026}

\begin{document}
\maketitle

{
\setcounter{tocdepth}{3}
\tableofcontents
\clearpage
}
Bản ngày 17/10/2026. Viết theo từng việc để người mới có thể đọc bản vẽ,
làm mô hình thử, đặt đúng phụ kiện và phối hợp với thợ lắp bản thật.

\textbf{Bắt đầu bằng một cửa mẫu và một tường xoay mẫu. Chỉ nhân ra toàn
bộ sau khi hai mẫu hoạt động ổn định và nhả khẩn được.} Bản 3D hiện tại
là thiết kế bố trí và logic trò chơi; chưa có tính toán kết cấu, bản vẽ
chế tạo trục, tải neo, khoảng chống kẹp hoặc thiết kế điện an toàn. Phần
vách cao 3 m, neo chịu lực và cơ cấu chuyển động dùng cho khách cần
người có chuyên môn chốt thiết kế trước khi làm. Hướng dẫn này không
biến các thông số mô phỏng thành kích thước cắt vật tư.

\section{1. Mở đúng bộ hồ
sơ}\label{mux1edf-ux111uxfang-bux1ed9-hux1ed3-sux1a1}

\begin{itemize}
\tightlist
\item
  \href{https://dam-truong.github.io/quy-mu-rev-p/model-360.html}{Viewer
  360°, bốn trạng thái và góc nhìn ra làng}.
\item
  \href{https://dam-truong.github.io/quy-mu-rev-p/blueprint/index.html}{Mặt
  bằng và các bản K-01 đến K-08}.
\item
  \href{https://dam-truong.github.io/quy-mu-rev-p/HUONG-DAN.html}{Báo
  cáo thay đổi REV O → REV P}.
\item
  \href{https://dam-truong.github.io/quy-mu-rev-p/src/layout.js}{Nguồn
  layout duy nhất}: dùng để đối chiếu vị trí, chiều cửa, trục xoay.
\item
  \href{https://dam-truong.github.io/quy-mu-rev-p/blueprint/K-08-nha-khan.html}{Phương
  án nhả khẩn E1} và
  \href{https://dam-truong.github.io/quy-mu-rev-p/emergency-plan.json}{kết
  quả kiểm tra}.
\item
  \href{https://dam-truong.github.io/quy-mu-rev-p/Quy-Mu-REV-P.zip}{Tải
  trọn bộ ZIP}.
\end{itemize}

REV P giữ thứ tự \textbf{Làng → Giam → Người canh → Tế → thoát}. Có 22
cửa một chiều, 4 tường xoay X1--X4, 2 khóa tuyến L1--L2 và 1 panel nhả
khẩn E1. C01--C06 là mã cửa ra/vào và cửa phòng; không cộng chúng vào số
22 cửa một chiều. Nhả khẩn là chức năng an toàn của cơ cấu hiện có,
không phải câu đố mới.

24 bot hoàn thành trong 616,46--729,51 giây ở 1,2 m/s, không cộng thời
gian chờ hay giải đố. Tuyến biết trước bản đồ chỉ khoảng 105,42 m, tương
đương 87,85 giây. Mốc trên 600 giây là kết quả phép thử bot, không phải
thời gian cam kết cho khách. Cả 8 lượt bám tường trái/phải bị lặp; không
được lấy lượt lặp làm bằng chứng khách sẽ chơi đủ lâu.

\section{2. Hiểu kích thước trước khi mua vật
tư}\label{hiux1ec3u-kuxedch-thux1b0ux1edbc-trux1b0ux1edbc-khi-mua-vux1eadt-tux1b0}

{\def\LTcaptype{none} % do not increment counter
\begin{longtable}[]{@{}
  >{\raggedright\arraybackslash}p{(\linewidth - 4\tabcolsep) * \real{0.2300}}
  >{\raggedright\arraybackslash}p{(\linewidth - 4\tabcolsep) * \real{0.1800}}
  >{\raggedright\arraybackslash}p{(\linewidth - 4\tabcolsep) * \real{0.5900}}@{}}
\toprule\noalign{}
\begin{minipage}[b]{\linewidth}\raggedright
Mục
\end{minipage} & \begin{minipage}[b]{\linewidth}\raggedright
Giá trị trong mô hình
\end{minipage} & \begin{minipage}[b]{\linewidth}\raggedright
Cách dùng khi thi công
\end{minipage} \\
\midrule\noalign{}
\endhead
\bottomrule\noalign{}
\endlastfoot
Khung bố trí & 11 × 11 m & Là khoảng giữa các đường biên mô hình; phải
xác nhận cách đặt tim/mặt vách với mặt bằng thật \\
Lưới & 12 × 12 ô & Đánh số x, z từ 0 đến 11 \\
Bước ô & 916,667 mm & Khoảng cách giữa hai đường lưới; không phải bề
rộng cánh cửa \\
Vách & dày 160 mm, cao 3.000 mm & Bao gồm phần hoàn thiện nếu cần giữ
đúng hình học \\
Cửa C01--C06 & tham số rộng 660 mm, cao 2.100 mm & Phải đo lại thông
thủy sau khung, bản lề, tay nắm và nẹp \\
Hành lang giữa hai vách song song cách một ô & khoảng 756,667 mm & Chỉ
là 916,667 − 160; chỗ giao, phụ kiện và cánh mở có thể làm hẹp hơn \\
\end{longtable}
}

Gốc ở góc trái trên mặt bằng; x tăng sang phải, z tăng xuống dưới. Tọa
độ \texttt{{[}x,z{]}} trong bảng cơ cấu là \textbf{ô}, còn tọa độ
trục/đoạn tường tính bằng \textbf{mét}. Tâm ô \texttt{{[}x,z{]}} là
\texttt{((x+0,5)\ ×\ 11/12;\ (z+0,5)\ ×\ 11/12)}.

Các đường \texttt{WALLS} là tim vách. Vách biên có bề dày nên bề mặt có
thể vượt đường tim 0 hoặc 11 m. Không dựng thẳng vào căn phòng thông
thủy 11 × 11 m rồi tự dịch vách vào trong: cách đó sẽ làm hẹp lối. Đo
hiện trường và thống nhất đường chuẩn trước; nếu phải đổi bố trí, sửa
layout và kiểm tra lại.

Chiều rộng trong bản này chưa phải xác nhận đáp ứng yêu cầu thoát nạn
hoặc tiếp cận của công trình. Người phụ trách công trình phải duyệt kích
thước sử dụng thực tế trước khi mở cho khách. Không thu hẹp lối để tăng
độ khó.

\section{3. Làm thử theo thứ tự
này}\label{luxe0m-thux1eed-theo-thux1ee9-tux1ef1-nuxe0y}

\begin{enumerate}
\def\labelenumi{\arabic{enumi}.}
\tightlist
\item
  \textbf{Mô hình bàn tỉ lệ 1:10:} nền 1.100 × 1.100 mm, ô 91,667 mm;
  dùng bìa/foam nhẹ. Vẽ tim bằng thước, dán mã vách/cửa. Làm bốn tường
  xoay bằng bìa và chốt nhỏ. Mô hình này chỉ kiểm tra bố trí, không kiểm
  tra tải trọng hay chống kẹp.
\item
  \textbf{Đánh dấu sàn thật:} căng hai đường vuông góc, kiểm tra hai
  đường chéo, chia lưới. Dán băng màu lên toàn bộ tim vách, cửa và vùng
  quét X1--X4. Đi thử Làng → Giam → Người canh → Tế → thoát khi chưa
  dựng vách.
\item
  \textbf{Dựng một cửa một chiều mẫu} bằng đúng khung và bộ phụ kiện dự
  kiến. Thử cả chế độ chơi và nhả khẩn.
\item
  \textbf{Dựng một tường xoay mẫu} trong khu thử có rào giới hạn. Đo
  hành trình, khe, độ võng và lực vận hành; dùng tải hoàn thiện thật sau
  khi hệ đỡ đã được duyệt.
\item
  \textbf{Chốt bản vẽ chế tạo:} ghi rõ kích thước từng khung/cánh, loại
  vật liệu, tải, phụ kiện, vị trí neo, khe lắp và vùng quét. Chỉ lúc này
  mới đặt cắt hàng loạt.
\item
  \textbf{Dựng vách cố định}, sau đó lắp khung cửa, cơ cấu, nhả khẩn,
  rồi mới ốp trang trí. Chừa nắp bảo trì nhìn thấy được từ phía nhân
  viên.
\item
  \textbf{Chạy nghiệm thu không có khách}, kiểm tra từng cơ cấu, từng
  pha và toàn bộ tuyến. Sau đó mới thử vận hành có giám sát.
\end{enumerate}

\section{4. Chuẩn bị đồ nghề và bảng vật
tư}\label{chuux1ea9n-bux1ecb-ux111ux1ed3-nghux1ec1-vuxe0-bux1ea3ng-vux1eadt-tux1b0}

Đồ đo/lắp: thước cuộn, thước vuông, máy cân bằng hoặc dây dọi, bút đánh
dấu, kẹp giữ, khoan theo vật liệu, dụng cụ siết đúng đầu, giá kê và
thiết bị nâng phù hợp. Cắt/khoan trên bàn làm việc; không giữ cánh nặng
bằng tay trong lúc bắt trục. Dùng bảo hộ theo công việc và hướng dẫn
dụng cụ.

{\def\LTcaptype{none} % do not increment counter
\begin{longtable}[]{@{}
  >{\raggedright\arraybackslash}p{(\linewidth - 4\tabcolsep) * \real{0.2300}}
  >{\raggedright\arraybackslash}p{(\linewidth - 4\tabcolsep) * \real{0.1800}}
  >{\raggedright\arraybackslash}p{(\linewidth - 4\tabcolsep) * \real{0.5900}}@{}}
\toprule\noalign{}
\begin{minipage}[b]{\linewidth}\raggedright
Cụm
\end{minipage} & \begin{minipage}[b]{\linewidth}\raggedright
Số bộ
\end{minipage} & \begin{minipage}[b]{\linewidth}\raggedright
Mỗi bộ cần có
\end{minipage} \\
\midrule\noalign{}
\endhead
\bottomrule\noalign{}
\endlastfoot
Cửa một chiều & 22 & Khung chịu lực; cánh nhẹ; bộ bản lề/trục phù hợp;
bộ tự đóng có chỉnh; cơ cấu hạn chế chiều có thể nhả; chặn hành trình;
che khe; tay đẩy/nhận diện; cách nhả bằng tay \\
Tường xoay & 4 & Khung cánh; hệ trục đỡ dưới và dẫn hướng trên đồng
trục; khung kết cấu đỡ; chặn ở hai vị trí; bộ giữ vị trí có nhả; che
vùng kẹp; tay vận hành ngoài vùng quét; xác nhận vị trí \\
Cửa ẩn L1, L2 & 2 & Khung; cánh; bộ ray/con trượt phù hợp; dẫn hướng;
chống nhấc/chống rơi; chặn cuối; giữ đóng/mở có nhả; tay mở khẩn; chỗ đỗ
cánh đã kiểm tra \\
Panel E1 & 1 & Khung; cánh được giữ chống đổ/rơi; cơ cấu mở nhanh bằng
tay; tay thao tác hai phía; giữ mở; biển nhận diện khi khẩn \\
Cửa C01--C06 & 6 vị trí & Bộ khung/cánh phù hợp chức năng, mở được theo
kế hoạch thoát nạn đã duyệt \\
Vách cố định & Theo bảng vách ở cuối & Khung, liên kết với kết cấu đã
duyệt, tấm ốp, nẹp bo cạnh, nắp bảo trì \\
\end{longtable}
}

Đây là số \textbf{cụm}, chưa phải số bản lề, vít hay mét thép. Không
chọn phụ kiện chỉ theo cân nặng: nhà cung cấp cần biết chiều cao, chiều
rộng, bề dày cánh, vị trí trọng tâm và tần suất sử dụng. Ví dụ tài liệu
bản lề lò xo Ives 3SP1 giới hạn theo cả kích thước cửa lẫn khối lượng;
bảng của mẫu đó không đủ để suy ra nó chịu được vách cao 3 m.
\href{https://us.allegion.com/content/dam/allegion-us-2/web-documents-2/InstallInstructions/Ives_3SPI_Spring_Hinge_Installation_Instructions_113960.pdf}{Hướng
dẫn chính hãng Ives}.

Điền phiếu này cho \textbf{mỗi loại cánh} trước khi đặt hàng:

{\def\LTcaptype{none} % do not increment counter
\begin{longtable}[]{@{}
  >{\raggedright\arraybackslash}p{(\linewidth - 2\tabcolsep) * \real{0.3000}}
  >{\raggedright\arraybackslash}p{(\linewidth - 2\tabcolsep) * \real{0.7000}}@{}}
\toprule\noalign{}
\begin{minipage}[b]{\linewidth}\raggedright
Thông tin
\end{minipage} & \begin{minipage}[b]{\linewidth}\raggedright
Giá trị phải điền
\end{minipage} \\
\midrule\noalign{}
\endhead
\bottomrule\noalign{}
\endlastfoot
Mã cánh / vị trí & \ldots{} \\
Kích thước phủ bì sau hoàn thiện & rộng \ldots{} / cao \ldots{} / dày
\ldots{} mm \\
Khối lượng đã cân, gồm ốp và phụ kiện & \ldots{} kg \\
Trọng tâm so với trục & \ldots{} mm; xác định bằng bản vẽ hoặc đo \\
Vật liệu và kết cấu nhận neo & \ldots{} \\
Bộ phụ kiện / tài liệu lắp / giới hạn áp dụng & \ldots{} \\
Khe quanh cánh và thông thủy nhỏ nhất & \ldots{} mm; có bản vẽ đối
chiếu \\
Lực vận hành đo được / tiêu chí được duyệt & \ldots{} N / \ldots{} \\
Người duyệt kết cấu và cơ khí / ngày & \ldots{} \\
\end{longtable}
}

\section{5. Dựng vách cố
định}\label{dux1ef1ng-vuxe1ch-cux1ed1-ux111ux1ecbnh}

\begin{enumerate}
\def\labelenumi{\arabic{enumi}.}
\tightlist
\item
  Đối chiếu bảng tim vách ở cuối với K-01. Đánh mã từng đoạn lên sàn và
  trên khung; nhận diện E1 để không bị ốp thành vách cố định.
\item
  Chọn hệ khung và tấm ốp phù hợp công trình. Lớp giả đá chỉ là trang
  trí, không chịu bản lề hoặc trục. Không dùng lớp trang trí giòn/nặng ở
  mép cánh chuyển động.
\item
  Gia công khung theo bản vẽ đã duyệt. Đo đường chéo để kiểm tra vuông;
  thử dựng, chống tạm trước khi liên kết. Vách cao 3 m cần hệ ổn định
  riêng, không chỉ đứng bằng chân nhỏ trên sàn.
\item
  Neo vào đúng kết cấu đã khảo sát. Không khoan neo theo khoảng cách tự
  đoán, không coi trần trang trí là điểm chịu lực. Loại neo, số lượng và
  chiều sâu do người thiết kế chọn theo nền thực tế.
\item
  Lắp thử khung cửa và cơ cấu khi chưa đóng kín hai mặt. Đo lại lối đi
  tại chỗ hẹp nhất, kể cả tay nắm và nẹp.
\item
  Chỉ ốp kín sau khi cửa hoạt động và còn tiếp cận được vít/trục/khóa để
  bảo trì. Bo các cạnh tiếp xúc; cố định dây và đồ trang trí ngoài vùng
  quét.
\end{enumerate}

\section{6. Cửa một chiều --- làm một bộ chuẩn rồi nhân 22
bộ}\label{cux1eeda-mux1ed9t-chiux1ec1u-luxe0m-mux1ed9t-bux1ed9-chuux1ea9n-rux1ed3i-nhuxe2n-22-bux1ed9}

\textbf{Cách hoạt động:} khách đi từ ô \texttt{from} sang ô \texttt{to}
thì đẩy qua được; cánh tự trở về. Ở chiều ngược lại, cơ cấu hạn chế
chiều giữ cửa trong chế độ chơi. Khi nhả khẩn, cả hai phía phải đi qua
được hoặc cánh được giữ hoàn toàn ngoài lối.

\begin{Shaded}
\begin{Highlighting}[]
\NormalTok{Chế độ chơi:       ô from  {-}{-}{-} đẩy qua {-}{-}{-}\textgreater{}  ô to}
\NormalTok{                   khung — trục — cánh — bộ hạn chế chiều}
\NormalTok{Nhả khẩn:          ô from  \textless{}{-}{-} đi qua {-}{-}{-}\textgreater{}  ô to}
\NormalTok{                   bộ hạn chế đã nhả; cánh không cản đường}
\end{Highlighting}
\end{Shaded}

\textbf{Chọn phần cứng trước:} yêu cầu nhà cung cấp một bộ có khả năng
nhả hạn chế chiều đúng yêu cầu trên. Bản lề chỉ mở một phía cộng nẹp
chặn cố định sẽ không tự trở thành cửa hai chiều khi cắt điện. Không làm
chốt tự chế nhả dưới tải mà chưa kiểm tra. Phương án có thể là cơ cấu
chặn có nhả và trục cho phép mở hai phía, hoặc giữ cánh ở vị trí đỗ hoàn
toàn thoáng; phải kiểm tra trên mẫu thật và trong mặt bằng.

Các bước lắp:

\begin{enumerate}
\def\labelenumi{\arabic{enumi}.}
\tightlist
\item
  Dán mũi tên \texttt{from\ →\ to} lên sàn. Chọn cạnh đặt trục và phía
  đỗ cánh bằng bản vẽ vùng quét; layout mới quy định \textbf{chiều đi},
  chưa quy định tay bản lề.
\item
  Cố định khung chịu lực, kiểm tra đứng và vuông. Đo ô khung hoàn thiện,
  rồi tính cánh theo khe của bộ phụ kiện; không cắt cánh rộng 916,667 mm
  chỉ vì đó là bước ô.
\item
  Đánh dấu lỗ theo dưỡng của nhà sản xuất. Kê cánh bằng giá; lắp đủ liên
  kết quy định. Với bản lề liên tục, hướng dẫn Ives yêu cầu lắp đầy đủ
  các vít được chỉ định; không thay bằng vài vít bắt vào tấm ốp.
  \href{https://us.allegion.com/content/dam/allegion-us-2/web-files/ives/installation-documents/Ives_700_715_705_711_Continuous_Hinge_Installation_Instructions_116335.pdf}{Tài
  liệu lắp Ives 700/715/705/711}.
\item
  Lắp tự đóng theo đúng tài liệu của bộ đã mua. Chỉnh mức thấp nhất vẫn
  đóng hết hành trình; không siết mạnh để làm trò chơi khó hơn. Kiểm tra
  cánh không quật vào người sau khi thả. Mỗi mẫu bản lề có giới hạn
  chỉnh riêng.
  \href{https://us.allegion.com/content/dam/allegion-us-2/web-documents-2/InstallInstructions/Ives_3SPI_Spring_Hinge_Installation_Instructions_113960.pdf}{Ví
  dụ quy trình của Ives 3SP1}.
\item
  Lắp hạn chế chiều, chặn hành trình và che khe trục. Chặn phải truyền
  lực về khung đủ khả năng chịu lực; không để lực đẩy cuối hành trình bẻ
  bản lề.
\item
  Lắp nhả bằng tay tại chỗ và chức năng nhả chung theo thiết kế an toàn.
  Nhả từ hai phía, kiểm tra khi đang có lực tì lên cánh bằng dụng cụ
  thử; không dùng người để tạo tình huống kẹt.
\item
  Đo lại lối đi khi mở hết. Nếu cánh/khung làm hẹp hành lang hoặc chặn
  cửa kế bên, dừng để sửa thiết kế, không đổi chiều mũi tên tùy ý.
\item
  Ghi mã OW01--OW22, hướng, loại phụ kiện, ngày lắp và kết quả thử. OW
  là mã thi công theo thứ tự trong \texttt{ONEWAY}; bảng tọa độ ở cuối
  được sinh từ cùng dữ liệu.
\end{enumerate}

\section{7. Tường xoay X1--X4 --- quay 90° giữa hai vị
trí}\label{tux1b0ux1eddng-xoay-x1x4-quay-90-giux1eefa-hai-vux1ecb-truxed}

Tường xoay là cánh vách đổi từ đoạn \texttt{before} sang đoạn
\texttt{after} quanh \texttt{post}. Trục là đường thẳng đứng; không phải
tâm cánh. X1/X2 đổi khi mở L1; X3/X4 đổi khi mở L2.

\begin{Shaded}
\begin{Highlighting}[]
\NormalTok{Nhìn ngang:                     Nhìn từ trên (ví dụ một cánh):}
\NormalTok{Khung kết cấu đỡ                 vị trí B}
\NormalTok{       | dẫn hướng trên             |}
\NormalTok{       |                            |  vùng quét 90°}
\NormalTok{       | cánh vách                  O{-}{-}{-}{-}{-}{-}{-}{-} vị trí A}
\NormalTok{       |                            O = trục}
\NormalTok{       O trục đỡ dưới}
\NormalTok{Kết cấu nền nhận tải}
\end{Highlighting}
\end{Shaded}

Quy trình làm mẫu thủ công:

\begin{enumerate}
\def\labelenumi{\arabic{enumi}.}
\tightlist
\item
  Từ bảng X ở cuối, đánh dấu trục và hai đường tim cánh trên sàn. Dùng
  dưỡng bìa bằng \textbf{kích thước phủ bì thực}, gồm nẹp và tay vận
  hành, xoay thử toàn bộ hành trình. Kiểm tra cả góc cánh, vùng sát
  trục, vách kế bên và đầu người. Hình 3D ở hai đầu hành trình chưa
  chứng minh phần cứng không va chạm ở giữa.
\item
  Chọn trọn bộ trục dưới chịu tải và dẫn hướng trên tương thích, cùng
  khung kết cấu nhận tải. Không treo tường cao 3 m bằng bản lề nội thất
  chọn theo cảm tính. Lắp các trục theo cùng một đường dọi; hệ pivot của
  Ives cũng yêu cầu phối hợp đúng bộ trục trên/dưới và phụ kiện trung
  gian.
  \href{https://us.allegion.com/content/dam/allegion-us-2/web-files/ives/installation-documents/Ives_7244F_7245F_1-1_2in_Intermediate_Pivot_Installation_Instructions_107133.pdf}{Tài
  liệu lắp pivot Ives}.
\item
  Làm cánh khung nhẹ, cứng, không vặn. Kích thước đầu trục, đầu tự do và
  khe lắp lấy từ thiết kế bộ pivot; đoạn tim dài một ô chưa phải chiều
  dài cắt cánh. Lắp bằng giá kê/thiết bị nâng; giữ chống đổ đến khi hệ
  đỡ hoàn chỉnh.
\item
  Lắp chặn cơ khí tại A và B, rồi bộ giữ vị trí. Không dùng chính động
  cơ hoặc ma sát trục làm chặn duy nhất. Cánh phải được giữ chắc khi
  khách tì vào ở mỗi vị trí.
\item
  Bố trí tay điều khiển phía nhân viên, ngoài vùng quét; che khe kẹp và
  phần truyền động. Không khoét lỗ tay vào vị trí có thể bị ép giữa cánh
  và vách.
\item
  Thử bằng tay trong khu đã cách ly. Đánh dấu điểm nặng/kẹt, kiểm tra độ
  nghiêng, võng, rung và va chạm. Không tăng lực để vượt chỗ kẹt; tìm
  nguyên nhân và sửa.
\item
  Thử đúng trình tự nhả E1 và 22 cửa một chiều trước khi chuyển vị trí.
  Sau khi cánh được giữ ở B mới xác nhận pha mới.
\end{enumerate}

\textbf{Bản có động cơ là một hạng mục riêng.} Cần thiết kế bảo vệ vùng
quét, phát hiện kẹt, giám sát vị trí/nhả, dừng khẩn và điều khiển an
toàn bởi người có chuyên môn. Không dùng Arduino, timer hay cảm biến
trang trí làm lớp bảo vệ duy nhất. Chưa có sơ đồ điện hoặc tải cơ khí
được duyệt thì chỉ thử mẫu thủ công trong vùng cách ly; không chạy tự
động cạnh khách. Thời gian hoạt cảnh trong viewer không phải tốc độ xoay
an toàn đã xác định.

\section{8. Cửa ẩn L1, L2 và cửa
phòng}\label{cux1eeda-ux1ea9n-l1-l2-vuxe0-cux1eeda-phuxf2ng}

L1 mở sau Giam để tiếp cận Người canh; L2 mở sau Người canh để tiếp cận
Tế. Chọn phương án mẫu \textbf{trượt bằng tay} với bộ ray đồng bộ có tải
phù hợp. Ray có thể đỡ dưới hoặc treo trên tùy thiết kế; nếu treo, kết
cấu đỡ phía trên phải được tính tải.

\begin{enumerate}
\def\labelenumi{\arabic{enumi}.}
\tightlist
\item
  Đánh dấu đoạn đóng từ bảng L. Vẽ cả cánh khi đóng, khi mở và đường
  trượt.
\item
  Xác nhận có \textbf{chỗ đỗ đủ cho toàn bộ cánh} và phụ kiện trong vùng
  cho phép. Layout hiện chưa thiết kế hốc giấu cánh. Không tự mở hốc vào
  phòng, vượt khung 11 × 11 m hoặc làm dày vách lấn hành lang. Nếu không
  có chỗ đỗ, cần thiết kế lại cơ cấu và kiểm tra layout trước khi chế
  tạo.
\item
  Lắp ray vào kết cấu chịu lực, cân thẳng theo hướng dẫn bộ ray; lắp
  chặn cuối trước khi thả con trượt chạy tự do.
\item
  Lắp chống nhấc/chống rơi và dẫn hướng; không bỏ dẫn hướng dưới chỉ để
  cửa nhìn đẹp. Lắp cánh rồi kiểm tra toàn hành trình không lệch hoặc tự
  trôi.
\item
  Lắp giữ đóng/giữ mở và tay nhả từ hai phía. Giấu đường ghép bằng hoa
  văn/sơn, không dùng nẹp bắc qua khe làm kẹt cửa.
\item
  Thử mở bằng tay khi mất điện và khi điều khiển trò chơi hỏng. Không
  dùng ổ khóa cần chìa riêng làm cách thoát duy nhất.
\end{enumerate}

C03--C05 giữ vị trí cửa phòng theo layout. C01 là lối vào; C02 là cổng
làng vào mê cung; C06 là cửa thoát trò chơi mở ở S3. Việc C06 bị khóa
trong mô phỏng \textbf{không cho phép khóa chặn thoát người trong công
trình thật}: cách mở khẩn và lối ra cuối cùng phải hoạt động độc lập với
tiến độ câu đố, được duyệt cùng phương án thoát nạn.

\section{9. E1 và cách nhả
khẩn}\label{e1-vuxe0-cuxe1ch-nhux1ea3-khux1ea9n}

E1 là panel trên đoạn \texttt{x\ =\ 2,75\ m}, từ \texttt{z\ =\ 6,416667}
đến \texttt{7,333333\ m}, nối ô \texttt{{[}2,7{]}\ ↔\ {[}3,7{]}}. Bình
thường đóng như vách. Làm cánh có hệ đỡ, mở/giữ mở được từ hai phía;
không dùng tấm nặng tháo rời phải nhấc ra mới đi qua. Vị trí bản lề và
chỗ đỗ cần được kiểm tra riêng để cánh mở không chiếm đường cứu hộ.

\textbf{Nhả khẩn phải giải phóng lối đi. Dừng động cơ thôi chưa đủ.} Khi
chuyển pha, kiểm tra đồ thị hiện tại giả định E1 mở và cả 22 cửa một
chiều cho đi hai chiều. Nếu bản thật không làm được điều đó, kết quả
kiểm tra đường về cổng không còn áp dụng.

Trình tự chuyển S0 → S1 hoặc S1 → S2:

\begin{enumerate}
\def\labelenumi{\arabic{enumi}.}
\tightlist
\item
  Dừng diễn/đuổi; xác nhận không có người hoặc vật trong toàn vùng
  chuyển động và vị trí đỗ cánh.
\item
  Mở E1, nhả hạn chế chiều của 22 cửa; xác nhận thực tế đã thông. Nhả
  bằng tay tuần tự chỉ dùng khi khu thử đã cách ly; vận hành có khách
  cần phương án nhả chung được thiết kế và kiểm chứng.
\item
  Giữ trạng thái nhả trong suốt chuyển động. Chỉ điều khiển cặp tường
  tương ứng khi điều kiện an toàn còn đúng; mất xác nhận thì dừng và giữ
  lối thoát.
\item
  Khi tường đến vị trí cuối, xác nhận chặn/giữ đã vào đúng chỗ rồi mới
  khôi phục trạng thái chơi.
\item
  Nếu có sự cố: dừng chuyển động, giữ các lối nhả, hướng người ra ngoài
  theo kế hoạch đã duyệt. Chỉ reset sau khi tìm và xử lý nguyên nhân;
  reset không được tự làm tường chạy lại. Nguyên tắc này phù hợp giải
  thích của Pilz về dừng khẩn: reset có chủ ý và không tự khởi động lại.
  \href{https://www.pilz.com/en-IE/support/faq/standards/articles/180045}{Pilz
  --- emergency stop reset}.
\end{enumerate}

Kiểm tra phần mềm đã xét trường hợp cả vị trí tường cũ và mới cùng bị
chặn, rồi tìm đường về C02. Nó chưa mô phỏng đầy đủ cơ thể người, vùng
quét liên tục, khói, tải đám đông, thời gian nhả hoặc tình huống linh
kiện hỏng. C02 cũng chưa phải bằng chứng đã đi ra nơi an toàn bên ngoài;
phải kiểm tra tiếp đường qua làng/C01 và lối thoát được duyệt.

\section{10. Bảng vận hành
S0--S3}\label{bux1ea3ng-vux1eadn-huxe0nh-s0s3}

{\def\LTcaptype{none} % do not increment counter
\begin{longtable}[]{@{}
  >{\raggedright\arraybackslash}p{(\linewidth - 12\tabcolsep) * \real{0.0600}}
  >{\raggedright\arraybackslash}p{(\linewidth - 12\tabcolsep) * \real{0.2200}}
  >{\raggedright\arraybackslash}p{(\linewidth - 12\tabcolsep) * \real{0.0800}}
  >{\raggedright\arraybackslash}p{(\linewidth - 12\tabcolsep) * \real{0.0800}}
  >{\raggedright\arraybackslash}p{(\linewidth - 12\tabcolsep) * \real{0.1600}}
  >{\raggedright\arraybackslash}p{(\linewidth - 12\tabcolsep) * \real{0.1600}}
  >{\raggedright\arraybackslash}p{(\linewidth - 12\tabcolsep) * \real{0.2400}}@{}}
\toprule\noalign{}
\begin{minipage}[b]{\linewidth}\raggedright
Pha
\end{minipage} & \begin{minipage}[b]{\linewidth}\raggedright
Điều kiện
\end{minipage} & \begin{minipage}[b]{\linewidth}\raggedright
L1
\end{minipage} & \begin{minipage}[b]{\linewidth}\raggedright
L2
\end{minipage} & \begin{minipage}[b]{\linewidth}\raggedright
X1, X2
\end{minipage} & \begin{minipage}[b]{\linewidth}\raggedright
X3, X4
\end{minipage} & \begin{minipage}[b]{\linewidth}\raggedright
C06 trong trò chơi
\end{minipage} \\
\midrule\noalign{}
\endhead
\bottomrule\noalign{}
\endlastfoot
S0 & Vừa ra làng & Đóng & Đóng & Before & Before & Đóng \\
S1 & Xong Giam & Mở & Đóng & After & Before & Đóng \\
S2 & Xong Người canh & Mở & Mở & After & After & Đóng \\
S3 & Xong Tế & Mở & Mở & After & After & Mở \\
\end{longtable}
}

Reset về S0 khi đã đưa hết người ra ngoài; áp dụng lại trình tự cách ly
và nhả trước khi xoay ngược. Không reset giữa lượt chơi chỉ để đưa cửa
về chỗ cũ.

\section{11. Nghiệm thu, bảo trì và xử lý
lỗi}\label{nghiux1ec7m-thu-bux1ea3o-truxec-vuxe0-xux1eed-luxfd-lux1ed7i}

In bảng này; ghi \textbf{Đạt / Không đạt / Người thử / Ngày}. Một mục
không đạt thì không đưa cụm liên quan vào phục vụ. Lịch kiểm tra cơ khí
và điện tuân theo nhà sản xuất và đánh giá sử dụng thực tế; số lần thử
tùy ý không thay thế nghiệm thu.

{\def\LTcaptype{none} % do not increment counter
\begin{longtable}[]{@{}
  >{\raggedright\arraybackslash}p{(\linewidth - 2\tabcolsep) * \real{0.3000}}
  >{\raggedright\arraybackslash}p{(\linewidth - 2\tabcolsep) * \real{0.7000}}@{}}
\toprule\noalign{}
\begin{minipage}[b]{\linewidth}\raggedright
Việc thử
\end{minipage} & \begin{minipage}[b]{\linewidth}\raggedright
Điều kiện đạt
\end{minipage} \\
\midrule\noalign{}
\endhead
\bottomrule\noalign{}
\endlastfoot
Kích thước và chiều cửa & Đúng bảng tọa độ; thông thủy thực tế và vùng
quét được duyệt; không thêm cơ cấu hoặc lấn phòng \\
Khung, neo, trục & Đã kiểm tra theo thiết kế; không lỏng, võng, lệch hay
cọ ở toàn hành trình \\
Từng OW01--OW22 & Đúng chiều khi chơi; tự đóng êm; nhả tại chỗ và nhả
chung đều giải phóng được lối \\
X1--X4 tại hai vị trí & Đúng trục/đoạn; giữ ổn định; không cản lối dự
kiến \\
X1--X4 giữa hành trình & Đã đo vùng quét thực; bảo vệ chống tiếp cận/kẹp
và phản ứng dừng được kiểm chứng bằng dụng cụ thử phù hợp \\
L1/L2 & Đúng thứ tự; không trượt khỏi ray; chỗ đỗ không chiếm phòng/lối;
mở khẩn được \\
E1 & Mở từ hai phía; giữ mở; không cần nhấc vật nặng, chìa hoặc tháo
vít \\
Mất điện / mất tín hiệu / nút khẩn & Không phát sinh chuyển động; vẫn mở
thoát được; reset không tự chạy lại \\
Đi toàn tuyến ở S0--S3 & Đúng trình tự, có đường về cổng; kiểm tra tiếp
đến lối ra an toàn thật \\
Biển và dấu dẫn & Nhìn được, không nằm sau cánh/vách hoặc nhầm với dấu
thoát nạn \\
Thử sau hoàn thiện & Sơn/nẹp/đạo cụ không làm tăng lực, gây kẹt, che
thao tác hoặc cản cảm biến \\
\end{longtable}
}

Trước mỗi ca: kiểm tra lối ra, E1, nhả một chiều, chốt giữ tường, vật lạ
trên ray và liên lạc với người vận hành. Nếu phát hiện lỗi thì dừng sử
dụng cụm lỗi, không buộc dây hoặc chêm tạm để tiếp tục chơi.

{\def\LTcaptype{none} % do not increment counter
\begin{longtable}[]{@{}
  >{\raggedright\arraybackslash}p{(\linewidth - 2\tabcolsep) * \real{0.3000}}
  >{\raggedright\arraybackslash}p{(\linewidth - 2\tabcolsep) * \real{0.7000}}@{}}
\toprule\noalign{}
\begin{minipage}[b]{\linewidth}\raggedright
Hiện tượng
\end{minipage} & \begin{minipage}[b]{\linewidth}\raggedright
Kiểm tra theo thứ tự
\end{minipage} \\
\midrule\noalign{}
\endhead
\bottomrule\noalign{}
\endlastfoot
Cửa một chiều không tự đóng & Vật cản → lệch khung/trục → cọ sàn/nẹp →
cài đặt tự đóng theo tài liệu \\
Cửa đóng quá mạnh & Dừng thử → kiểm tra bộ tự đóng/giảm chấn → chỉnh
theo tài liệu → đo lại lực \\
Tường xoay nặng một đoạn & Dừng → kiểm tra đồng trục, nền, võng cánh và
va chạm; không tăng công suất để ép qua \\
Cửa trượt bị nghiêng & Cách ly → kê/giữ cánh an toàn → kiểm tra ray, con
trượt, chống nhấc và dẫn hướng \\
Nhả nhưng vẫn không qua được & Kiểm tra chặn cố định, tay bản lề, vị trí
đỗ và cửa kế bên; sửa cơ khí rồi thử lại \\
Mất điện còn khóa người & Không vận hành; sửa phương án nhả độc lập
trước khi thử lại \\
\end{longtable}
}

Sau mọi thay đổi về vách, chiều cửa, vị trí trục, vùng đỗ hoặc bề rộng
lối: cập nhật \texttt{src/layout.js}/nguồn sinh tương ứng, chạy lại
pipeline kiểm tra rồi nghiệm thu thực tế phần bị ảnh hưởng. Thay phụ
kiện cùng chức năng vẫn phải kiểm tra lại khe, lực, tải và nhả khẩn.

\section{12. Bảng tọa độ lấy trực tiếp từ REV
P}\label{bux1ea3ng-tux1ecda-ux111ux1ed9-lux1ea5y-trux1ef1c-tiux1ebfp-tux1eeb-rev-p}

Các bảng dưới được sinh tự động từ layout đã kiểm tra. Giữ nguyên mã để
trao đổi với thợ. Số làm tròn phục vụ đọc; dùng nguồn layout khi triển
khai đo đạc. \textbf{Bảng tim vách không phải bảng cắt:} đầu nối, bề
dày, trục, khung, khe và E1 phải được triển khai thành chi tiết trước
khi đặt vật tư.

SHA-256 layout:
\nolinkurl{4bfa818db193497c01c28394ab6103fea1207aca01d534ef65c586cb47f82a6b}.

\subsection{12.1. Cửa ra/vào và cửa phòng
(mét)}\label{cux1eeda-ravuxe0o-vuxe0-cux1eeda-phuxf2ng-muxe9t}

{\def\LTcaptype{none} % do not increment counter
\begin{longtable}[]{@{}
  >{\raggedright\arraybackslash}p{(\linewidth - 6\tabcolsep) * \real{0.1000}}
  >{\raggedright\arraybackslash}p{(\linewidth - 6\tabcolsep) * \real{0.2900}}
  >{\raggedright\arraybackslash}p{(\linewidth - 6\tabcolsep) * \real{0.2200}}
  >{\raggedright\arraybackslash}p{(\linewidth - 6\tabcolsep) * \real{0.3900}}@{}}
\toprule\noalign{}
\begin{minipage}[b]{\linewidth}\raggedright
Mã
\end{minipage} & \begin{minipage}[b]{\linewidth}\raggedright
Tâm x; z (m)
\end{minipage} & \begin{minipage}[b]{\linewidth}\raggedright
Trục đoạn cửa
\end{minipage} & \begin{minipage}[b]{\linewidth}\raggedright
Chức năng
\end{minipage} \\
\midrule\noalign{}
\endhead
\bottomrule\noalign{}
\endlastfoot
C01 & (0.0000; 9.6250) & z & Vào làng \\
C02 & (4.1250; 8.2500) & x & Làng → mê cung \\
C03 & (10.0833; 0.4583) & z & Giam \\
C04 & (10.5417; 8.2500) & x & Người canh \\
C05 & (1.8333; 2.2917) & z & Tế \\
C06 & (11.0000; 0.4583) & z & Thoát trò chơi \\
\end{longtable}
}

Trục x: đoạn cửa nằm ngang mặt bằng; trục z: đoạn cửa nằm dọc. Đây không
phải ký hiệu tay bản lề.

\subsection{12.2. Cửa một chiều}\label{cux1eeda-mux1ed9t-chiux1ec1u}

Đoạn đặt cửa là ranh giới chung của hai ô. Bảng ghi đoạn tim, chưa trừ
khung/khe.

{\def\LTcaptype{none} % do not increment counter
\begin{longtable}[]{@{}
  >{\raggedright\arraybackslash}p{(\linewidth - 6\tabcolsep) * \real{0.1000}}
  >{\raggedright\arraybackslash}p{(\linewidth - 6\tabcolsep) * \real{0.1400}}
  >{\raggedright\arraybackslash}p{(\linewidth - 6\tabcolsep) * \real{0.1400}}
  >{\raggedright\arraybackslash}p{(\linewidth - 6\tabcolsep) * \real{0.6200}}@{}}
\toprule\noalign{}
\begin{minipage}[b]{\linewidth}\raggedright
Mã
\end{minipage} & \begin{minipage}[b]{\linewidth}\raggedright
Đi từ ô
\end{minipage} & \begin{minipage}[b]{\linewidth}\raggedright
Sang ô
\end{minipage} & \begin{minipage}[b]{\linewidth}\raggedright
Đoạn ranh giới (m)
\end{minipage} \\
\midrule\noalign{}
\endhead
\bottomrule\noalign{}
\endlastfoot
OW01 & {[}2, 5{]} & {[}2, 4{]} & (1.8333; 4.5833) → (2.7500; 4.5833) \\
OW02 & {[}3, 5{]} & {[}3, 4{]} & (2.7500; 4.5833) → (3.6667; 4.5833) \\
OW03 & {[}8, 4{]} & {[}7, 4{]} & (7.3333; 3.6667) → (7.3333; 4.5833) \\
OW04 & {[}8, 5{]} & {[}8, 4{]} & (7.3333; 4.5833) → (8.2500; 4.5833) \\
OW05 & {[}2, 6{]} & {[}2, 5{]} & (1.8333; 5.5000) → (2.7500; 5.5000) \\
OW06 & {[}3, 6{]} & {[}3, 5{]} & (2.7500; 5.5000) → (3.6667; 5.5000) \\
OW07 & {[}4, 5{]} & {[}5, 5{]} & (4.5833; 4.5833) → (4.5833; 5.5000) \\
OW08 & {[}4, 6{]} & {[}4, 5{]} & (3.6667; 5.5000) → (4.5833; 5.5000) \\
OW09 & {[}5, 5{]} & {[}6, 5{]} & (5.5000; 4.5833) → (5.5000; 5.5000) \\
OW10 & {[}9, 6{]} & {[}9, 5{]} & (8.2500; 5.5000) → (9.1667; 5.5000) \\
OW11 & {[}1, 7{]} & {[}1, 6{]} & (0.9167; 6.4167) → (1.8333; 6.4167) \\
OW12 & {[}5, 6{]} & {[}6, 6{]} & (5.5000; 5.5000) → (5.5000; 6.4167) \\
OW13 & {[}5, 7{]} & {[}5, 6{]} & (4.5833; 6.4167) → (5.5000; 6.4167) \\
OW14 & {[}7, 7{]} & {[}7, 6{]} & (6.4167; 6.4167) → (7.3333; 6.4167) \\
OW15 & {[}9, 7{]} & {[}9, 6{]} & (8.2500; 6.4167) → (9.1667; 6.4167) \\
OW16 & {[}2, 7{]} & {[}1, 7{]} & (1.8333; 6.4167) → (1.8333; 7.3333) \\
OW17 & {[}2, 8{]} & {[}2, 7{]} & (1.8333; 7.3333) → (2.7500; 7.3333) \\
OW18 & {[}4, 7{]} & {[}5, 7{]} & (4.5833; 6.4167) → (4.5833; 7.3333) \\
OW19 & {[}5, 7{]} & {[}5, 8{]} & (4.5833; 7.3333) → (5.5000; 7.3333) \\
OW20 & {[}6, 7{]} & {[}7, 7{]} & (6.4167; 6.4167) → (6.4167; 7.3333) \\
OW21 & {[}6, 7{]} & {[}6, 8{]} & (5.5000; 7.3333) → (6.4167; 7.3333) \\
OW22 & {[}5, 8{]} & {[}4, 8{]} & (4.5833; 7.3333) → (4.5833; 8.2500) \\
\end{longtable}
}

\subsection{12.3. Tường xoay}\label{tux1b0ux1eddng-xoay}

{\def\LTcaptype{none} % do not increment counter
\begin{longtable}[]{@{}
  >{\raggedright\arraybackslash}p{(\linewidth - 8\tabcolsep) * \real{0.0700}}
  >{\raggedright\arraybackslash}p{(\linewidth - 8\tabcolsep) * \real{0.1000}}
  >{\raggedright\arraybackslash}p{(\linewidth - 8\tabcolsep) * \real{0.1900}}
  >{\raggedright\arraybackslash}p{(\linewidth - 8\tabcolsep) * \real{0.3200}}
  >{\raggedright\arraybackslash}p{(\linewidth - 8\tabcolsep) * \real{0.3200}}@{}}
\toprule\noalign{}
\begin{minipage}[b]{\linewidth}\raggedright
Mã
\end{minipage} & \begin{minipage}[b]{\linewidth}\raggedright
Khi mở
\end{minipage} & \begin{minipage}[b]{\linewidth}\raggedright
Trục x; z (m)
\end{minipage} & \begin{minipage}[b]{\linewidth}\raggedright
Before --- đoạn tim (m)
\end{minipage} & \begin{minipage}[b]{\linewidth}\raggedright
After --- đoạn tim (m)
\end{minipage} \\
\midrule\noalign{}
\endhead
\bottomrule\noalign{}
\endlastfoot
X1 & L1 & (10.0833; 6.4167) & (9.1667; 6.4167) → (10.0833; 6.4167) &
(10.0833; 5.5000) → (10.0833; 6.4167) \\
X2 & L1 & (8.2500; 6.4167) & (8.2500; 6.4167) → (8.2500; 7.3333) &
(7.3333; 6.4167) → (8.2500; 6.4167) \\
X3 & L2 & (0.9167; 2.7500) & (0.9167; 2.7500) → (1.8333; 2.7500) &
(0.9167; 2.7500) → (0.9167; 3.6667) \\
X4 & L2 & (3.6667; 7.3333) & (3.6667; 6.4167) → (3.6667; 7.3333) &
(2.7500; 7.3333) → (3.6667; 7.3333) \\
\end{longtable}
}

X4 của REV P đặt trục tại (3,6667; 7,3333) m; không lấy hướng X4 từ REV
O để lắp. Chọn cung quay 90° nối hai vị trí rồi kiểm tra vùng quét bằng
cánh thật; không quay vòng dài 270°.

\subsection{12.4. Khóa tuyến và panel
E1}\label{khuxf3a-tuyux1ebfn-vuxe0-panel-e1}

{\def\LTcaptype{none} % do not increment counter
\begin{longtable}[]{@{}
  >{\raggedright\arraybackslash}p{(\linewidth - 6\tabcolsep) * \real{0.0800}}
  >{\raggedright\arraybackslash}p{(\linewidth - 6\tabcolsep) * \real{0.1900}}
  >{\raggedright\arraybackslash}p{(\linewidth - 6\tabcolsep) * \real{0.4600}}
  >{\raggedright\arraybackslash}p{(\linewidth - 6\tabcolsep) * \real{0.2700}}@{}}
\toprule\noalign{}
\begin{minipage}[b]{\linewidth}\raggedright
Mã
\end{minipage} & \begin{minipage}[b]{\linewidth}\raggedright
Hai ô
\end{minipage} & \begin{minipage}[b]{\linewidth}\raggedright
Đoạn tim (m)
\end{minipage} & \begin{minipage}[b]{\linewidth}\raggedright
Mở khi
\end{minipage} \\
\midrule\noalign{}
\endhead
\bottomrule\noalign{}
\endlastfoot
L1 & {[}10, 8{]} ↔ {[}11, 8{]} & (10.0833; 7.3333) → (10.0833; 8.2500) &
Xong Giam \\
L2 & {[}1, 4{]} ↔ {[}1, 5{]} & (0.9167; 4.5833) → (1.8333; 4.5833) &
Xong Người canh \\
E1 & {[}2, 7{]} ↔ {[}3, 7{]} & (2.7500; 6.4167) → (2.7500; 7.3333) & Nhả
khẩn / trước mọi chuyển pha có xoay \\
\end{longtable}
}

\subsection{12.5. Tim các đoạn
vách}\label{tim-cuxe1c-ux111oux1ea1n-vuxe1ch}

Mã W dưới đây dùng để đánh dấu thi công theo thứ tự \texttt{WALLS}. Một
đoạn dài có thể chia nhiều module theo thiết kế khung. \textbf{Đoạn W có
E1 phải được thay bằng panel có thể mở, không xây kín.} Các trục
xoay/cửa được liệt kê riêng, không cộng nhầm vào vách cố định.

{\def\LTcaptype{none} % do not increment counter
\begin{longtable}[]{@{}
  >{\raggedright\arraybackslash}p{(\linewidth - 6\tabcolsep) * \real{0.0800}}
  >{\raggedright\arraybackslash}p{(\linewidth - 6\tabcolsep) * \real{0.4600}}
  >{\raggedright\arraybackslash}p{(\linewidth - 6\tabcolsep) * \real{0.1500}}
  >{\raggedright\arraybackslash}p{(\linewidth - 6\tabcolsep) * \real{0.3100}}@{}}
\toprule\noalign{}
\begin{minipage}[b]{\linewidth}\raggedright
Mã
\end{minipage} & \begin{minipage}[b]{\linewidth}\raggedright
Tim đầu → cuối (m)
\end{minipage} & \begin{minipage}[b]{\linewidth}\raggedright
Chiều dài tim (m)
\end{minipage} & \begin{minipage}[b]{\linewidth}\raggedright
Ghi chú
\end{minipage} \\
\midrule\noalign{}
\endhead
\bottomrule\noalign{}
\endlastfoot
W01 & (0.0000; 0.0000) → (11.0000; 0.0000) & 11.0000 & Triển khai đầu
nối trước khi cắt \\
W02 & (0.9167; 0.9167) → (1.8333; 0.9167) & 0.9166 & Triển khai đầu nối
trước khi cắt \\
W03 & (0.9167; 1.8333) → (1.8333; 1.8333) & 0.9166 & Triển khai đầu nối
trước khi cắt \\
W04 & (6.4167; 1.8333) → (10.0833; 1.8333) & 3.6666 & Triển khai đầu nối
trước khi cắt \\
W05 & (1.8333; 2.7500) → (6.4167; 2.7500) & 4.5834 & Triển khai đầu nối
trước khi cắt \\
W06 & (9.1667; 2.7500) → (10.0833; 2.7500) & 0.9166 & Triển khai đầu nối
trước khi cắt \\
W07 & (3.6667; 3.6667) → (4.5833; 3.6667) & 0.9166 & Triển khai đầu nối
trước khi cắt \\
W08 & (6.4167; 3.6667) → (9.1667; 3.6667) & 2.7500 & Triển khai đầu nối
trước khi cắt \\
W09 & (4.5833; 4.5833) → (6.4167; 4.5833) & 1.8334 & Triển khai đầu nối
trước khi cắt \\
W10 & (9.1667; 4.5833) → (10.0833; 4.5833) & 0.9166 & Triển khai đầu nối
trước khi cắt \\
W11 & (0.0000; 5.5000) → (0.9167; 5.5000) & 0.9167 & Triển khai đầu nối
trước khi cắt \\
W12 & (4.5833; 5.5000) → (5.5000; 5.5000) & 0.9167 & Triển khai đầu nối
trước khi cắt \\
W13 & (6.4167; 5.5000) → (7.3333; 5.5000) & 0.9166 & Triển khai đầu nối
trước khi cắt \\
W14 & (9.1667; 5.5000) → (10.0833; 5.5000) & 0.9166 & Triển khai đầu nối
trước khi cắt \\
W15 & (2.7500; 6.4167) → (3.6667; 6.4167) & 0.9167 & Triển khai đầu nối
trước khi cắt \\
W16 & (5.5000; 6.4167) → (6.4167; 6.4167) & 0.9167 & Triển khai đầu nối
trước khi cắt \\
W17 & (10.0833; 6.4167) → (11.0000; 6.4167) & 0.9167 & Triển khai đầu
nối trước khi cắt \\
W18 & (0.9167; 7.3333) → (1.8333; 7.3333) & 0.9166 & Triển khai đầu nối
trước khi cắt \\
W19 & (3.6667; 7.3333) → (4.5833; 7.3333) & 0.9166 & Triển khai đầu nối
trước khi cắt \\
W20 & (6.4167; 7.3333) → (7.3333; 7.3333) & 0.9166 & Triển khai đầu nối
trước khi cắt \\
W21 & (8.2500; 7.3333) → (10.0833; 7.3333) & 1.8333 & Triển khai đầu nối
trước khi cắt \\
W22 & (0.0000; 8.2500) → (3.6667; 8.2500) & 3.6667 & Triển khai đầu nối
trước khi cắt \\
W23 & (4.5833; 8.2500) → (10.0833; 8.2500) & 5.5000 & Triển khai đầu nối
trước khi cắt \\
W24 & (0.0000; 11.0000) → (11.0000; 11.0000) & 11.0000 & Triển khai đầu
nối trước khi cắt \\
W25 & (0.0000; 0.0000) → (0.0000; 9.1667) & 9.1667 & Triển khai đầu nối
trước khi cắt \\
W26 & (0.0000; 10.0833) → (0.0000; 11.0000) & 0.9167 & Triển khai đầu
nối trước khi cắt \\
W27 & (0.9167; 1.8333) → (0.9167; 2.7500) & 0.9167 & Triển khai đầu nối
trước khi cắt \\
W28 & (0.9167; 3.6667) → (0.9167; 4.5833) & 0.9166 & Triển khai đầu nối
trước khi cắt \\
W29 & (0.9167; 6.4167) → (0.9167; 7.3333) & 0.9166 & Triển khai đầu nối
trước khi cắt \\
W30 & (1.8333; 0.0000) → (1.8333; 1.8333) & 1.8333 & Triển khai đầu nối
trước khi cắt \\
W31 & (1.8333; 2.7500) → (1.8333; 3.6667) & 0.9167 & Triển khai đầu nối
trước khi cắt \\
W32 & (1.8333; 4.5833) → (1.8333; 6.4167) & 1.8334 & Triển khai đầu nối
trước khi cắt \\
W33 & (2.7500; 3.6667) → (2.7500; 5.5000) & 1.8333 & Triển khai đầu nối
trước khi cắt \\
W34 & (2.7500; 6.4167) → (2.7500; 7.3333) & 0.9166 & E1 --- panel nhả,
bình thường đóng \\
W35 & (3.6667; 3.6667) → (3.6667; 5.5000) & 1.8333 & Triển khai đầu nối
trước khi cắt \\
W36 & (4.5833; 3.6667) → (4.5833; 4.5833) & 0.9166 & Triển khai đầu nối
trước khi cắt \\
W37 & (4.5833; 5.5000) → (4.5833; 6.4167) & 0.9167 & Triển khai đầu nối
trước khi cắt \\
W38 & (5.5000; 2.7500) → (5.5000; 3.6667) & 0.9167 & Triển khai đầu nối
trước khi cắt \\
W39 & (5.5000; 7.3333) → (5.5000; 8.2500) & 0.9167 & Triển khai đầu nối
trước khi cắt \\
W40 & (6.4167; 0.0000) → (6.4167; 2.7500) & 2.7500 & Triển khai đầu nối
trước khi cắt \\
W41 & (7.3333; 1.8333) → (7.3333; 3.6667) & 1.8334 & Triển khai đầu nối
trước khi cắt \\
W42 & (7.3333; 4.5833) → (7.3333; 7.3333) & 2.7500 & Triển khai đầu nối
trước khi cắt \\
W43 & (8.2500; 2.7500) → (8.2500; 4.5833) & 1.8333 & Triển khai đầu nối
trước khi cắt \\
W44 & (8.2500; 5.5000) → (8.2500; 6.4167) & 0.9167 & Triển khai đầu nối
trước khi cắt \\
W45 & (8.2500; 8.2500) → (8.2500; 11.0000) & 2.7500 & Triển khai đầu nối
trước khi cắt \\
W46 & (9.1667; 1.8333) → (9.1667; 2.7500) & 0.9167 & Triển khai đầu nối
trước khi cắt \\
W47 & (9.1667; 6.4167) → (9.1667; 7.3333) & 0.9166 & Triển khai đầu nối
trước khi cắt \\
W48 & (10.0833; 0.9167) → (10.0833; 1.8333) & 0.9166 & Triển khai đầu
nối trước khi cắt \\
W49 & (10.0833; 2.7500) → (10.0833; 4.5833) & 1.8333 & Triển khai đầu
nối trước khi cắt \\
W50 & (10.0833; 6.4167) → (10.0833; 7.3333) & 0.9166 & Triển khai đầu
nối trước khi cắt \\
W51 & (11.0000; 0.9167) → (11.0000; 11.0000) & 10.0833 & Triển khai đầu
nối trước khi cắt \\
\end{longtable}
}

\subsection{12.6. Phiếu bàn giao từng
cụm}\label{phiux1ebfu-buxe0n-giao-tux1eebng-cux1ee5m}

\begin{itemize}
\tightlist
\item
  Mã cụm: \ldots; người lắp: \ldots; ngày: \ldots{}
\item
  Bản vẽ chế tạo và tài liệu phụ kiện đính kèm: \ldots{}
\item
  Kích thước/khối lượng thực tế: \ldots{}
\item
  Chiều mở/vùng quét/vị trí đỗ đã đối chiếu: \ldots{}
\item
  Thử chế độ chơi: \ldots; thử nhả khẩn/mất điện: \ldots{}
\item
  Lỗi còn lại và cách xử lý: \ldots{}
\item
  Người kiểm tra chấp nhận cho sử dụng: \ldots{}
\end{itemize}

Giữ phiếu cạnh hồ sơ bảo trì. Nếu thay đổi phần cứng, cập nhật phiếu
thay vì dựa vào ảnh chụp bản cũ.

\end{document}
